\documentclass[11pt]{article}
\usepackage{amsmath, amssymb}
\usepackage[margin=1in]{geometry}

\title{Derivation of the Diagonal-Covariance Gaussian\\Log-Probability for MDN Loss}
\author{keras-mdn-layer}
\date{}

\begin{document}
\maketitle

\section*{Starting Point: Bishop PRML Eq.~2.43}

The multivariate Gaussian distribution in $D$ dimensions is:
\begin{equation}
    \mathcal{N}(\mathbf{y} \mid \boldsymbol{\mu}, \boldsymbol{\Sigma})
    = \frac{1}{(2\pi)^{D/2}\,|\boldsymbol{\Sigma}|^{1/2}}
      \exp\!\left\{
        -\frac{1}{2}(\mathbf{y}-\boldsymbol{\mu})^\top
        \boldsymbol{\Sigma}^{-1}
        (\mathbf{y}-\boldsymbol{\mu})
      \right\}
    \tag{PRML 2.43}
\end{equation}

\section*{Diagonal Covariance Assumption}

In a mixture density network we parameterise each component with a
diagonal covariance matrix:
\begin{equation}
    \boldsymbol{\Sigma}
    = \operatorname{diag}(\sigma_1^2,\, \sigma_2^2,\, \dots,\, \sigma_D^2)
\end{equation}

Three quantities simplify immediately:

\paragraph{1.\ Determinant.}
\begin{equation}
    |\boldsymbol{\Sigma}|
    = \prod_{d=1}^{D} \sigma_d^2
    \qquad\Longrightarrow\qquad
    |\boldsymbol{\Sigma}|^{1/2}
    = \prod_{d=1}^{D} \sigma_d
\end{equation}

\paragraph{2.\ Inverse.}
\begin{equation}
    \boldsymbol{\Sigma}^{-1}
    = \operatorname{diag}\!\left(
        \frac{1}{\sigma_1^2},\, \dots,\, \frac{1}{\sigma_D^2}
      \right)
\end{equation}

\paragraph{3.\ Quadratic form.}
Because $\boldsymbol{\Sigma}^{-1}$ is diagonal, the matrix product
reduces to an element-wise sum:
\begin{equation}
    (\mathbf{y}-\boldsymbol{\mu})^\top
    \boldsymbol{\Sigma}^{-1}
    (\mathbf{y}-\boldsymbol{\mu})
    = \sum_{d=1}^{D}
      \frac{(y_d - \mu_d)^2}{\sigma_d^2}
\end{equation}

\section*{Substitution}

Substituting (2)--(4) into (PRML 2.43):
\begin{equation}
    \mathcal{N}(\mathbf{y} \mid \boldsymbol{\mu}, \boldsymbol{\Sigma})
    = \frac{1}{(2\pi)^{D/2}\;\prod_{d}\sigma_d}
      \;\exp\!\left\{
        -\frac{1}{2}\sum_{d=1}^{D}
        \left(\frac{y_d - \mu_d}{\sigma_d}\right)^{\!2}
      \right\}
\end{equation}

\section*{Taking the Logarithm}

\begin{equation}
\boxed{
    \log \mathcal{N}(\mathbf{y} \mid \boldsymbol{\mu}, \boldsymbol{\Sigma})
    = \underbrace{-\frac{D}{2}\log(2\pi)}_{\text{normalisation}}
      \;-\; \underbrace{\sum_{d=1}^{D}\log\sigma_d}_{\text{volume}}
      \;-\; \underbrace{\frac{1}{2}\sum_{d=1}^{D}
            \left(\frac{y_d - \mu_d}{\sigma_d}\right)^{\!2}}_{\text{squared Mahalanobis}}
}
\end{equation}

\section*{Correspondence to Code}

Each of the three terms maps directly to a line in
\texttt{keras\_mdn\_layer/\_\_init\_\_.py}:

\begin{verbatim}
log_component = (
    -0.5 * output_dim * math.log(2.0 * math.pi)              # -D/2 log(2pi)
    - ops.sum(ops.log(sigma), axis=-1)                        # -sum_d log(sigma_d)
    - 0.5 * ops.sum(ops.square((y_true - mu) / sigma), axis=-1)  # -1/2 sum_d (...)^2
)
\end{verbatim}

\end{document}
